% Chapter 1 — Notation and conventions. This chapter is the authoritative
% statement of decision D-7 and of every notational rule the remaining
% chapters, the source code, the log format, and the APIs must follow.
\chapter{Notation and Conventions}
\label{ch:notation}

\section{Scope and authority}

This chapter is normative for the entire project. Every quantity that appears
in this document, in the source code, in the SRLOG binary log, and in any
public API follows the conventions defined here. Where an external reference
uses a different convention, the transcription into this project's convention
is stated explicitly at the point of use. The attitude conventions in
Section~\ref{sec:notation:quaternion} ratify project decision D-7.

\section{Reference frames}
\label{sec:notation:frames}

Frames are denoted by single calligraphic letters used as sub- and
superscripts on vectors and direction-cosine matrices. The frames used across
the project are listed in Table~\ref{tab:notation:frames}. In Phase~1 only the
GCRF is exercised; the remaining frames are named now so that channel and API
labels are stable, and each is defined rigorously in the chapter of the phase
that introduces it.

\begin{table}[htbp]
  \centering
  \caption{Project reference frames.}
  \label{tab:notation:frames}
  \begin{tabular}{clp{7.2cm}}
    \toprule
    Symbol & Frame & Definition and status \\
    \midrule
    $I$ & GCRF &
      Geocentric Celestial Reference Frame: Earth-centered, kinematically
      non-rotating, axes aligned with the International Celestial Reference
      System per the IERS Conventions (2010)~\cite{iers2010}. The Phase~1
      integration and logging frame. \\
    $E$ & Earth-fixed &
      Rotating terrestrial frame reached from the GCRF through the CIO-based
      transformation chain of IERS Conventions (2010) Chapter~5. Defined in
      the Phase~2 frames chapter. \\
    $L$ & Moon principal-axis &
      Lunar body-fixed principal-axis frame realized from the DE440 libration
      angles~\cite{park2021}. Defined in the Phase~2 frames chapter. \\
    $M$ & Mars-fixed &
      Mars body-fixed frame from the IAU 2015 rotational elements. Defined in
      the Phase~2 frames chapter, where its defining reference is cited. \\
    $B$ & Body &
      Vehicle structural frame, $+X$ forward. Defined in the Phase~4 vehicle
      chapter. In Phase~1 the body frame exists only as the target frame of
      the logged placeholder attitude quaternion. \\
    \bottomrule
  \end{tabular}
\end{table}

\section{Vectors, matrices, and the DCM convention}
\label{sec:notation:dcm}

Physical vectors are set in bold italic, $\vv{r}$, and matrices in bold
upright-italic capitals, $\mat{C}$. The coordinates of a vector resolved in
frame $A$ are written $\vv{r}^{A}$. The skew-symmetric cross-product matrix of
$\vv{a} = [a_1, a_2, a_3]^{\mathsf T}$ is
\begin{equation}
  \skewm{\vv{a}} =
  \begin{bmatrix}
    0 & -a_3 & a_2 \\
    a_3 & 0 & -a_1 \\
    -a_2 & a_1 & 0
  \end{bmatrix},
  \qquad
  \skewm{\vv{a}}\,\vv{b} = \vv{a} \times \vv{b}.
  \label{eq:notation:skew}
\end{equation}

A direction-cosine matrix (DCM) written $\dcm{A}{B}$ \emph{maps frame $A$ to
frame $B$}; that is, it transforms the coordinates of a vector from frame $A$
to frame $B$:
\begin{equation}
  \vv{v}^{B} = \dcm{A}{B}\,\vv{v}^{A}.
  \label{eq:notation:dcmdef}
\end{equation}
DCMs are orthonormal, so the inverse transformation is the transpose,
$\bigl(\dcm{A}{B}\bigr)^{-1} = \bigl(\dcm{A}{B}\bigr)^{\mathsf T} =
\dcm{B}{A}$, and successive transformations compose by left multiplication:
\begin{equation}
  \dcm{A}{C} = \dcm{B}{C}\,\dcm{A}{B}.
  \label{eq:notation:dcmcomp}
\end{equation}

\section{Quaternion convention (D-7)}
\label{sec:notation:quaternion}

\subsection{Hamilton quaternions, scalar-first}

The project uses \emph{Hamilton} quaternions: the basis elements satisfy
$i^2 = j^2 = k^2 = ijk = -1$, and hence $ij = k$. Writing a quaternion as a
scalar part $q_w$ and a vector part $\vv{q}_v = [q_x, q_y, q_z]^{\mathsf T}$,
the Hamilton product of $\quat{p}$ and $\quat{q}$ is
\begin{equation}
  \quat{p} \quatmul \quat{q} =
  \begin{bmatrix}
    p_w q_w - \vv{p}_v \cdot \vv{q}_v \\
    p_w \vv{q}_v + q_w \vv{p}_v + \vv{p}_v \times \vv{q}_v
  \end{bmatrix}.
  \label{eq:notation:hamiltonproduct}
\end{equation}
The alternative (often called JPL) convention, which flips the sign of the
cross-product term, is \emph{not} used anywhere in this project; the
distinction between the two conventions is surveyed by Markley and
Crassidis~\cite{markley2014}.

Quaternion components are ordered \textbf{scalar-first} in \emph{all}
documentation, logs, and APIs:
\begin{equation}
  \quat{q} = [\,q_w,\; q_x,\; q_y,\; q_z\,],
  \qquad
  \text{identity} = [\,1,\; 0,\; 0,\; 0\,].
  \label{eq:notation:scalarfirst}
\end{equation}
Attitude quaternions are unit quaternions, $\norm{\quat{q}} = 1$, and the
inverse of a unit quaternion is its conjugate,
$\quat{q}^{-1} = \quat{q}^{*} = [\,q_w, -q_x, -q_y, -q_z\,]$.

\subsection{The attitude quaternion \texorpdfstring{$\quat{q}_{i2b}$}{q\_i2b}}

The logged and exchanged attitude quaternion is $\quat{q}_{i2b}$, the frame
transformation \emph{from the inertial frame $I$ (GCRF) to the body frame
$B$}. Embedding a vector as the pure quaternion
$\tilde{\vv{v}} = [\,0, \vv{v}\,]$, coordinates transform by conjugation:
\begin{equation}
  \tilde{\vv{v}}^{B}
    = \quat{q}_{i2b}^{-1} \quatmul \tilde{\vv{v}}^{I} \quatmul \quat{q}_{i2b}.
  \label{eq:notation:sandwich}
\end{equation}
Equivalently, the DCM corresponding to $\quat{q}_{i2b}$ is
\begin{equation}
  \dcm{I}{B}(\quat{q})
    = \bigl(q_w^2 - \vv{q}_v \cdot \vv{q}_v\bigr)\,\mat{I}_3
      + 2\,\vv{q}_v \vv{q}_v^{\mathsf T}
      - 2\,q_w \skewm{\vv{q}_v},
  \label{eq:notation:quat2dcm}
\end{equation}
which is the attitude-matrix form given by Markley and
Crassidis~\cite{markley2014}, transcribed here to the scalar-first component
ordering of equation~\eqref{eq:notation:scalarfirst}. Under
equations~\eqref{eq:notation:hamiltonproduct}
and~\eqref{eq:notation:sandwich}, frame transformations compose in
left-to-right chaining order,
\begin{equation}
  \quat{q}_{i2c} = \quat{q}_{i2b} \quatmul \quat{q}_{b2c},
  \qquad
  \dcm{I}{B}(\quat{p} \quatmul \quat{q})
    = \dcm{I}{B}(\quat{q})\,\dcm{I}{B}(\quat{p}),
  \label{eq:notation:quatcomp}
\end{equation}
consistent with the DCM composition rule of
equation~\eqref{eq:notation:dcmcomp}. Quaternion kinematics and rigid-body
attitude dynamics are specified in the attitude-dynamics chapter that lands
with the vehicle phase; Phase~1 logs only the fixed identity quaternion.

\subsection{Mandatory note: mapping to \texttt{Eigen::Quaterniond}}
\label{sec:notation:eigen}

The C++ core represents quaternions with
\texttt{Eigen::Quaterniond}~\cite{eigenweb}, whose component ordering differs
between its constructor and its internal storage. This mismatch is the single
most likely source of attitude-convention defects, so the mapping is spelled
out here and is binding on all core and binding code:

\begin{itemize}
  \item The value constructor takes the scalar \emph{first}:
    \texttt{Eigen::Quaterniond(w, x, y, z)}.
  \item The internal coefficient storage, exposed by \texttt{.coeffs()},
    orders the scalar \emph{last}: \texttt{[x, y, z, w]}.
\end{itemize}

\begin{table}[htbp]
  \centering
  \caption{Component mapping between the project convention and
    \texttt{Eigen::Quaterniond}.}
  \label{tab:notation:eigenmap}
  \begin{tabular}{lccc}
    \toprule
    Component & Project / SRLOG / API index & Eigen accessor &
      \texttt{coeffs()} index \\
    \midrule
    $q_w$ (scalar) & 0 & \texttt{.w()} & 3 \\
    $q_x$          & 1 & \texttt{.x()} & 0 \\
    $q_y$          & 2 & \texttt{.y()} & 1 \\
    $q_z$          & 3 & \texttt{.z()} & 2 \\
    \bottomrule
  \end{tabular}
\end{table}

Consequences, stated as rules:
\begin{enumerate}
  \item Serialization code must never copy \texttt{.coeffs()} verbatim into a
    log record or API buffer; every externally visible quaternion is emitted
    scalar-first via the named accessors, per
    Table~\ref{tab:notation:eigenmap}.
  \item Eigen's \texttt{operator*} on vectors and
    \texttt{.toRotationMatrix()} implement the \emph{active} rotation
    $\vv{v} \mapsto \quat{q} \quatmul \tilde{\vv{v}} \quatmul
    \quat{q}^{-1}$. Applied to $\quat{q}_{i2b}$, therefore,
    \texttt{.toRotationMatrix()} returns $\dcm{B}{I}$ --- the \emph{transpose}
    of the frame transformation in
    equation~\eqref{eq:notation:quat2dcm}. Code that needs $\dcm{I}{B}$
    conjugates first or transposes the result, and says so in a comment.
\end{enumerate}

\section{Units}
\label{sec:notation:units}

All quantities are SI throughout the core, the logs, and the APIs; no
non-SI unit crosses any interface. In this document, numerical values with
units are typeset with the \texttt{siunitx} package, for example
$\mu_\oplus = \qty{3.986004418e14}{\metre\cubed\per\second\squared}$, the
geocentric gravitational constant (TCB-compatible value) of the IERS
Conventions (2010)~\cite{iers2010}. Configuration keys and log channel names
carry their unit as a suffix (\texttt{thrust\_vac\_N}, \texttt{r\_m},
\texttt{v\_mps}, \texttt{w\_b\_radps}, \texttt{mass\_kg}, \texttt{t\_s}) so
that a wrong unit is visible at the point of use. Angles are radians in code
and logs; degrees appear only in human-facing configuration keys explicitly
suffixed \texttt{\_deg}.

Epochs in configuration files are UTC ISO-8601 strings. The internal
timescale (TAI, decision D-6) and all timescale conversions are defined in
the Phase~2 time-systems chapter.

\section{Symbols}
\label{sec:notation:symbols}

Table~\ref{tab:notation:symbols} lists the symbols used project-wide.
Chapter-local symbols are defined where they first appear.

\begin{table}[htbp]
  \centering
  \caption{Project-wide symbols.}
  \label{tab:notation:symbols}
  \begin{tabular}{clc}
    \toprule
    Symbol & Meaning & Unit \\
    \midrule
    $t$ & simulation time since the mission epoch & \unit{\second} \\
    $\vv{r}$ & position vector & \unit{\metre} \\
    $\vv{v}$ & velocity vector & \unit{\metre\per\second} \\
    $\vv{a}$ & acceleration vector & \unit{\metre\per\second\squared} \\
    $\quat{q}_{i2b}$ & attitude quaternion, GCRF to body & 1 \\
    $\vv{\omega}^{B}$ & angular velocity of $B$ w.r.t.\ $I$, in $B$ & \unit{\radian\per\second} \\
    $m$ & spacecraft mass & \unit{\kilogram} \\
    $G$ & universal gravitational constant & \unit{\metre\cubed\per\kilogram\per\second\squared} \\
    $\mu$ & gravitational parameter $GM$ of the central body & \unit{\metre\cubed\per\second\squared} \\
    $\dcm{A}{B}$ & DCM mapping frame $A$ to frame $B$ & 1 \\
    $\mat{I}_3$ & $3 \times 3$ identity matrix & 1 \\
    $\vv{h}$ & specific angular momentum $\vv{r} \times \vv{v}$ & \unit{\metre\squared\per\second} \\
    $\varepsilon$ & specific orbital energy & \unit{\metre\squared\per\second\squared} \\
    $a$ & semi-major axis & \unit{\metre} \\
    $T$ & orbital period & \unit{\second} \\
    \bottomrule
  \end{tabular}
\end{table}
